\documentclass[conference]{IEEEtran}
\usepackage{cite}
\usepackage{amsmath,amssymb}
\usepackage{booktabs}
\usepackage{array}
\usepackage{url}
\usepackage[T1]{fontenc}
\newcolumntype{L}[1]{>{\raggedright\arraybackslash}p{#1}}

\title{How Do Migratory Animals Know Where They Are Going?\\
A Design-Only Framework for Testing Multisensory Navigation and Magnetoreception}
\author{Survey--Idea--ExperimentDesign--Author Research Workflow}

\begin{document}
\maketitle

\begin{abstract}
Migration is not one sensory act. A long-distance migrant must decide when to leave, establish a direction, compensate for the changing position of the sun and stars, respond to wind or water flow, recognize familiar places, estimate displacement from a target region, and select safe habitat. These tasks occur at different spatial and temporal scales and cannot be inferred from a single successful track. Magnetic information is particularly informative because the geomagnetic field can provide a directional compass and, in principle, spatial gradients. However, a behavioural bearing shift, a molecular magnetic-field effect, and an ecologically useful map are different claims. This design-only paper introduces MAGNAV, the Multisensory Animal GeoNavigation Assessment and Validation framework, for testing those claims without treating one cue or receptor as a universal explanation of migration. MAGNAV separates cue availability, cue-dependent orientation, receptor or neural mediation, route-level contribution, and cross-context generalization. It requires a species-, season-, developmental-stage-, and route-specific boundary; preregistered cue-conflict and restoration controls; a cue-dependence ledger; and noncompensatory gates for manipulation integrity, welfare, competing cues, and replication. The framework uses an S0--S4 scenario ladder spanning sensory-context atlases, controlled orientation tests, receptor-to-behaviour bridge assays, free-ranging validation, and independent replication. It treats cryptochrome 4 magnetic photochemistry in European robin as important discovery-stage evidence for a radical-pair candidate, not proof that the protein is the in vivo receptor or that magnetic sensing fully explains migration. No animal, molecular, tracking, or clinical-style outcome is reported here. The conclusion is conditional: migratory navigation can be studied as multiscale cue integration, but credible claims require separate evidence for the particular species, cue, computation, ecological context, and mechanism under discussion.
\end{abstract}

\begin{IEEEkeywords}
animal migration, navigation, magnetoreception, cryptochrome, radical pairs, celestial compass, cue integration, movement ecology.
\end{IEEEkeywords}

\section{Introduction}
The everyday phrase ``animals know where they are going'' compresses several biological problems into one. An animal can maintain a southwesterly bearing without knowing its geographic position. It can recognize a familiar coastline without possessing a global map. It can respond to a magnetic field in a molecular assay without using that response during a migratory flight. These distinctions matter because migration covers insects, birds, fish, turtles, marine mammals, and other taxa operating in radically different sensory environments.

Long-distance routes are especially demanding. A migrant must generate an initial heading, retain or update it while cues move or disappear, detect perturbations such as wind displacement, and make route or stopover decisions. Mouritsen's synthesis emphasizes that no single cue can provide pinpoint accuracy over thousands of kilometres; multiscale and multisensory integration is expected \cite{mouritsen2018}. The appropriate scientific question is therefore conditional: which information source is used by a particular animal, at a defined migratory phase and environmental context, and what navigation computation does it support?

The prompt foregrounds magnetoreception. Magnetic cues are biologically plausible and experimentally tractable, but they must be interpreted with care. A light-dependent inclination compass, magnetic-map-like responses, magnetite hypotheses, and induction-based sensitivity are not interchangeable. The radical-pair hypothesis is strengthened by in vitro magnetic sensitivity of European robin cryptochrome 4 (ErCRY4) \cite{xu2021}; it does not establish the complete receptor-to-neural-circuit-to-route chain. This paper converts that boundary into a staged research architecture called MAGNAV.

\section{Survey: Why One Navigation Explanation Fails}\label{sec:survey}
\subsection{Navigation is layered}
Three layers should be separated. A compass supplies direction relative to a reference frame. A map supplies a position-dependent relationship between the animal and a goal. A route controller combines these estimates with energetic state, weather, habitat, social information, and time. Celestial, visual, olfactory, hydrodynamic, and geomagnetic information may contribute at more than one layer, but evidence for one layer does not automatically transfer to another.

For individual $i$ at time $t$, a measured horizontal bearing $\theta_{it}$ can be represented by the unit vector
\begin{equation}
\mathbf{u}_{it}=\left(\cos\theta_{it},\sin\theta_{it}\right),
\label{eq:unitvector}
\end{equation}
where the coordinate origin and angle convention are frozen before analysis. Equation~\eqref{eq:unitvector} records a direction, not knowledge of a destination. For a release block with weights $w_i$, the mean resultant vector is
\begin{equation}
\bar{\mathbf{u}}_t=\frac{\sum_i w_i\mathbf{u}_{it}}{\sum_i w_i},\qquad R_t=\lVert\bar{\mathbf{u}}_t\rVert .
\label{eq:meanresultant}
\end{equation}
Here $R_t$ is concentration, so a shift in the mean can be distinguished from loss of orientation. Equation~\eqref{eq:meanresultant} prevents a diffuse or bimodal response from being misreported as a single preferred heading.

\begin{table*}[!t]
\caption{Evidence boundary for claims about migratory navigation.}\label{tab:boundary}\centering\footnotesize
\begin{tabular}{@{}L{.19\textwidth}L{.25\textwidth}L{.27\textwidth}L{.21\textwidth}@{}}\toprule
Evidence class & What it can establish & Principal limitations & Unsupported shortcut\\\midrule
Orientation arena response & Conditional change in bearing or concentration under stated cues. & Arena, light, handling, motivation, and cue availability may differ from migration. & A geographic map or wild-route consequence.\\
Celestial cue manipulation & Time- or polarization-dependent directional information. & Clock, weather, retinal state, and learned calibration remain possible mediators. & A universal solar or stellar mechanism.\\
Magnetic virtual displacement & Sensitivity to a stated field vector or map-like gradient. & Coil artefacts, electromagnetic noise, and retained cues can confound interpretation. & Identification of the receptor or route-level navigation.\\
Molecular magnetic effect & Candidate photochemical or cellular mechanism under assay conditions. & Expression, localization, signalling, and in vivo causal role remain untested. & Proof of animal-level magnetoreception.\\
Free-ranging track & Ecological movement outcome in a sampled season and location. & Wind, survival, tag effects, route experience, and missingness are coupled to movement. & A single-cue causal explanation.\\
Replication & Reproducibility within a locked boundary. & Finite species, seasons, and conditions remain finite. & A mechanism shared by all migrants.\\\bottomrule
\end{tabular}
\end{table*}

\subsection{Celestial, landscape, and magnetic information}
Solar and stellar patterns can provide directional reference frames, while visual landscapes, odours, currents, and waves can add local or regional information. A solar compass requires time compensation because solar azimuth changes during the day; monarch antennae have been shown to participate in the circadian organization of sun-compass navigation \cite{merlin2009}, and polarized skylight can remain informative when direct solar observations are degraded \cite{muheim2006}. For a clock phase $\tau$ and observed solar azimuth $\theta_{\mathrm{sun}}(t)$, a generic corrected compass bearing can be expressed as
\begin{equation}
\theta_{\mathrm{celestial}}(t)=\theta_{\mathrm{sun}}(t)-g(\tau,t,\ell),
\label{eq:timecomp}
\end{equation}
where $g$ is a species- and location-dependent calibration function and $\ell$ denotes latitude or relevant geometry. Equation~\eqref{eq:timecomp} is a model of the required computation, not a claim that every migratory species implements the same function.

The geomagnetic field contains declination, inclination, and intensity. These may serve as directional or spatial cues depending on the animal and task; sea-turtle work illustrates why a magnetic compass and an open-sea navigation account must still be evaluated separately \cite{lohmann1996}. The radical-pair account predicts a light-dependent chemical response to field direction and strength \cite{ritz2000,hore2016}. Consistent with this class of mechanism, Xu \emph{et al.} reported that ErCRY4 photochemistry is magnetically sensitive in vitro and characterized radical-pair pathways that can produce magnetic-field effects \cite{xu2021}. Birds also have behavioural evidence that visual rather than trigeminal mediation can carry magnetic-compass information in a defined migratory context \cite{zapka2009}. None of these observations makes beak magnetite, retinal cryptochrome, and all magnetic behaviours the same causal system.

\begin{table*}[!t]
\caption{Cue-to-computation hypotheses that must be tested separately.}\label{tab:cue-computation}\centering\footnotesize
\begin{tabular}{@{}L{.18\textwidth}L{.20\textwidth}L{.28\textwidth}L{.26\textwidth}@{}}\toprule
Cue family & Candidate computation & Discriminating manipulation & Bounded conclusion\\\midrule
Sun, stars, polarized light & Directional compass with time or sky calibration. & Clock shift, sky-pattern rotation, light-spectrum control, cue restoration. & Celestial information contributes to orientation under the tested condition.\\
Landmarks, olfaction, habitat & Local recognition or regional map calibration. & Familiar versus unfamiliar release, odour masking, panorama control, return test. & Specified landscape or odour information affects local navigation.\\
Magnetic inclination/declination/intensity & Compass, virtual displacement, or map-like gradient cue. & Controlled field rotation or displacement with physical-field measurements and sham coils. & Behaviour depends on the stated magnetic manipulation.\\
Radical-pair candidate & Light-dependent receptor or transduction step. & Wavelength/RF prediction, receptor variant or pathway perturbation, matched sensory controls. & Bounded receptor-to-response contribution.\\
Electromagnetic induction & Motion-through-water signal in electroreceptive animals. & Flow, conductivity, and field-vector manipulation with electroreception controls. & Induction-like information is available and behaviourally relevant.\\\bottomrule
\end{tabular}
\end{table*}

\section{Idea: MAGNAV}\label{sec:idea}
MAGNAV denotes \emph{Multisensory Animal GeoNavigation Assessment and Validation}. It is not an algorithm assigned to an animal. It is a claim-control framework that records what cue was manipulated, what cues remained available, which computation was tested, and which inference is licensed. It therefore avoids reducing a navigation system to the cue that happened to be easiest to manipulate.

For cue set $C_{it}$, environmental state $X_{it}$, and internal state $Z_{it}$, a generic integrated directional decision can be represented as
\begin{equation}
\hat{\mathbf{d}}_{it}=\frac{\sum_{c\in C_{it}}\alpha_c(X_{it},Z_{it})\mathbf{v}_{c,it}}{\left\lVert\sum_{c\in C_{it}}\alpha_c(X_{it},Z_{it})\mathbf{v}_{c,it}\right\rVert},
\label{eq:integration}
\end{equation}
where $\mathbf{v}_{c,it}$ is a cue-derived directional vector and $\alpha_c$ is an adaptive weight. Equation~\eqref{eq:integration} captures a testable possibility: a magnetic cue may dominate when celestial information is unavailable but be downweighted when sky cues or landmarks are reliable. It does not identify the weights without data.

The framework advances through five evidence levels: cue availability, cue-dependent orientation, mechanism-to-behaviour mediation, ecological route contribution, and independent extension. A result stays at its level until the next gate has been passed. MAGNAV's central object is a \emph{cue-dependence ledger}, which records the denominator of sampled animals, retained cues, target manipulation, restoration control, exclusions, and strongest permitted conclusion.

\begin{table*}[!t]
\caption{MAGNAV scenario ladder.}\label{tab:ladder}\centering\footnotesize
\begin{tabular}{@{}L{.11\textwidth}L{.25\textwidth}L{.26\textwidth}L{.29\textwidth}@{}}\toprule
Stage & System and intervention & Primary measurement & Advancement condition\\\midrule
S0 & Species, season, developmental state, route, light, weather, and magnetic-context atlas. & Provenance, sensor coverage, field and sky metadata. & Sampling and exposure boundaries are locked.\\
S1 & Controlled orientation with cue rotation, virtual displacement, cross-cue conflict, sham, and restoration. & Bearing, concentration, latency, and individual heterogeneity. & Predicted contrast appears against sham and reverses or calibrates as declared.\\
S2 & Candidate receptor, sensory pathway, or neural computation test with matched controls. & Manipulation integrity, cue response, sensory function, and viability. & Target perturbation is interpretable and mechanism prediction survives controls.\\
S3 & Free-ranging, ethically tagged or observational study with track-quality and weather provenance. & Route deviation, correction latency, stopover decision, and missingness. & Ecological effect is distinguishable from weather, device, or release effects.\\
S4 & Independent laboratory, site, season, or specified taxon with frozen analysis. & Replication, restriction, and negative-result reporting. & The context boundary replicates or is honestly narrowed.\\\bottomrule
\end{tabular}
\end{table*}

\section{ExperimentDesign: Controlled Tests to Ecological Consequence}\label{sec:design}
\subsection{Population boundary and primary estimands}
The initial demonstrator should be one preregistered nocturnally migrating songbird population in one season, not ``birds'' as a generic category. S0 fixes age/experience class, capture or observation geography, calendar window, moon and cloud conditions, light spectrum, local field vector, anthropogenic electromagnetic environment, and route history where available. Such detail is not administrative overhead; it specifies which cue combinations the experiment can possibly distinguish.

In S1, the primary estimand is the change in circular bearing between a stated cue-conflict arm and a matched sham arm. Let $A_i\in\{0,1\}$ indicate assigned arm and $Y_i$ be a preregistered angular or vector response. The contrast is
\begin{equation}
\Delta_{\mathrm{bearing}}=\mathbb{E}\left[Y_i\mid A_i=1,B_i=1\right]-\mathbb{E}\left[Y_i\mid A_i=0,B_i=1\right],
\label{eq:bearingcontrast}
\end{equation}
where $B_i=1$ denotes eligibility after predeclared quality checks. Equation~\eqref{eq:bearingcontrast} must be estimated jointly with concentration and the complete eligibility ledger. It is not meaningful if the experimental and sham field, illumination, or handling conditions are physically distinguishable in another way.

S3 evaluates route-level consequences rather than an arena proxy. For route segment $s$, a map-relevant directional error can be defined as
\begin{equation}
E_{is}=\operatorname{ang}\left(\mathbf{h}_{is},\mathbf{g}_{is}\right),
\label{eq:routeerror}
\end{equation}
where $\mathbf{h}_{is}$ is observed displacement heading and $\mathbf{g}_{is}$ is a preregistered goal or expected corridor vector. Equation~\eqref{eq:routeerror} requires that $\mathbf{g}_{is}$ be defined without looking at the manipulated tracks. A low error alone does not identify which retained cue was used.

\begin{table*}[!t]
\caption{Endpoint hierarchy and prohibited inference.}\label{tab:endpoints}\centering\footnotesize
\begin{tabular}{@{}L{.19\textwidth}L{.25\textwidth}L{.25\textwidth}L{.23\textwidth}@{}}\toprule
Endpoint tier & Required measurement & Interpretation enabled & Prohibited shortcut\\\midrule
Manipulation integrity & Calibrated magnetic vector, light spectrum, RF noise, release and handling records. & Whether an arm differs in the intended physical cue. & Assuming a coil or light device altered only one cue.\\
Orientation response & Individual bearing, concentration, latency, sham, and restoration. & Conditional cue-dependent orientation. & Calling a bearing shift a map or receptor result.\\
Mechanism bridge & Candidate-pathway perturbation, non-target control, sensory-function and health checks. & Bounded causal contribution under the assay. & Claiming complete in vivo magnetoreception.\\
Ecological route outcome & Track error, correction timing, weather/current, tag and data-loss provenance. & Route-level contribution in the sampled context. & Generalizing one season to a species-wide map.\\
Replication & Locked protocol, independent samples, full negative and attrition results. & Reproducibility of a defined claim. & A universal navigation mechanism.\\\bottomrule
\end{tabular}
\end{table*}

\subsection{Cue conflicts and restoration}
The decisive intervention is not simple cue removal. MAGNAV combines baseline, magnetic rotation or virtual displacement, celestial-clock or polarization manipulation where ethical and feasible, cross-cue conflict, sham, and restoration conditions. The planned magnetic rotation is represented by
\begin{equation}
\mathbf{B}^{\prime}_{it}=\mathbf{Q}(\phi)\mathbf{B}_{it},
\label{eq:fieldrotation}
\end{equation}
where $\mathbf{B}_{it}$ is the measured baseline vector and $\mathbf{Q}(\phi)$ is a verified rotation by angle $\phi$. Equation~\eqref{eq:fieldrotation} formalizes a physical intervention; it does not guarantee that the animal experiences no accompanying auditory, thermal, visual, or electromagnetic artefact.

A restoration test asks whether the response returns toward the preregistered baseline relation when the conflict is removed. The effect-size decomposition
\begin{equation}
\mathcal{C}=\Delta_{\mathrm{conflict}}-\Delta_{\mathrm{sham}},\qquad \mathcal{R}=\Delta_{\mathrm{restore}}-\Delta_{\mathrm{sham}},
\label{eq:restore}
\end{equation}
separates a planned conflict contrast $\mathcal{C}$ from restoration contrast $\mathcal{R}$. Equation~\eqref{eq:restore} must be interpreted with circular-appropriate estimators and uncertainty, but it prevents a one-way perturbation from being called cue-specific without a return condition.

\begin{table*}[!t]
\caption{Confound and provenance ledger.}\label{tab:confounds}\centering\footnotesize
\begin{tabular}{@{}L{.20\textwidth}L{.27\textwidth}L{.23\textwidth}L{.22\textwidth}@{}}\toprule
Threat & Required record or control & Failure if omitted & Resulting limit\\\midrule
Coil or device artefact & Field mapping, sham hardware, thermal/acoustic/light monitoring. & Behaviour reacts to apparatus rather than magnetic vector. & No cue-specific conclusion.\\
Celestial confounding & Spectrum, irradiance, cloud, polarization proxy, lunar phase, clock phase. & Magnetic and celestial reliability covary. & No magnetic-versus-celestial weighting inference.\\
Wind, current, and release geometry & High-resolution flow/weather and exact release metadata. & Route deviation reflects transport or site bias. & No map correction claim.\\
Health, device burden, and handling & Welfare thresholds, mass/drag, duration, recovery, stop rules. & Disorientation reflects stress or impairment. & No sensory-mechanism claim.\\
Track missingness & Battery/duty cycle, positional error, loss reason, censoring time. & Selective attrition changes observed route distribution. & No ecological effect estimate without sensitivity analysis.\\
Experience and social context & Age, prior migration, group state, and release block. & Cue use varies with learning or following. & No population-wide generalization.\\\bottomrule
\end{tabular}
\end{table*}

\section{Mechanism-to-Behaviour Bridge}\label{sec:mechanism}
The radical-pair mechanism should be treated as a bridge hypothesis. Light absorption by a cryptochrome can generate spin-correlated radical pairs whose chemical yields vary with magnetic conditions. The relevant inference is conditional: if a candidate protein has an appropriate photochemical response, localization, expression, signalling connection, and perturbation-to-behaviour association, it becomes a stronger mechanistic explanation for that population and condition. It is still not a species-general conclusion.

For a molecular or neural readout $M$, a target perturbation $P$, and cue condition $A$, a bridge interaction may be summarized as
\begin{align}
\Gamma={}&\left[\mathbb{E}(M\mid P=1,A=1)-\mathbb{E}(M\mid P=1,A=0)\right]\nonumber\\
&-\left[\mathbb{E}(M\mid P=0,A=1)-\mathbb{E}(M\mid P=0,A=0)\right].
\label{eq:interaction}
\end{align}
Equation~\eqref{eq:interaction} asks whether cue responsiveness depends on the nominated perturbation. It cannot repair non-specific perturbation, altered vision, toxicity, developmental compensation, or a mismatch between the assay and migratory state.

\begin{table*}[!t]
\caption{Mechanism claims and the evidence needed before advancing them.}\label{tab:mechanism}\centering\footnotesize
\begin{tabular}{@{}L{.20\textwidth}L{.24\textwidth}L{.28\textwidth}L{.20\textwidth}@{}}\toprule
Candidate statement & Necessary evidence & Critical alternatives to exclude & Maximum claim\\\midrule
CRY4-like radical pair is magnetically responsive & Reproduced field-dependent chemistry with calibrated light and molecular controls. & Photochemical artefact, non-physiological concentration, unrelated signalling. & Molecular candidate.\\
Candidate receptor mediates a compass response & Localization and target-specific perturbation paired with sensory-function and rescue controls. & Vision or health impairment, off-target effects, developmental compensation. & Bounded receptor-to-orientation bridge.\\
Magnetite-like particles contribute & Independent localization, material identification, perturbation, and behavioural rescue logic. & Contamination, macrophage/iron storage, non-specific lesion effects. & Context-specific magnetic-sensing contribution.\\
Induction supports aquatic orientation & Flow, conductivity, field, and electroreception manipulations with matched hydrodynamic controls. & Current-following, pressure, vibration, or non-magnetic electroreception. & Conditional induction-like cue use.\\\bottomrule
\end{tabular}
\end{table*}

\section{Field Translation, Analysis, and Reproducibility}\label{sec:translation}
Laboratory orientation and field navigation answer different questions. An arena can expose cue conflicts precisely but has limited ecological realism. Tracking has ecological realism but introduces weather, social, device, and data-loss confounding. A valid programme links them through predeclared predictions rather than assuming that a positive laboratory result will appear unchanged in the field.

Track inclusion should not erase missingness. For retained follow-up indicator $R_i(t)$ and covariate history $L_i(t)$, an inverse-retention sensitivity weight can be written as
\begin{equation}
W_i(t)=\prod_{s\leq t}\frac{q_s(1\mid A_i)}{q_s(1\mid A_i,L_i(s))},
\label{eq:retention}
\end{equation}
where $q_s$ models the probability that usable tracking information remains at time $s$. Equation~\eqref{eq:retention} exposes assumptions about selective loss; it does not create information about animals whose outcome is unobserved.

The experiment also needs a model for heterogeneity. A cue may be useful only under cloud cover, at a particular latitude, before learning a route, or during a transition between day and night. Therefore, a cue-by-context interaction is a planned scientific result, not a failed universal story. The multiplicity family should be frozen before data collection. For $m$ ordered confirmatory hypotheses, a step-down rule may compare
\begin{equation}
p_{(j)}\leq\frac{\alpha}{m-j+1},
\label{eq:holm}
\end{equation}
at rank $j$. Equation~\eqref{eq:holm} controls a stated family-wise error target but cannot correct post hoc selection of species, route, cue definition, or outcome window.

\begin{table*}[!t]
\caption{Decision gates and value-of-information priorities.}\label{tab:gates}\centering\footnotesize
\begin{tabular}{@{}L{.18\textwidth}L{.27\textwidth}L{.27\textwidth}L{.20\textwidth}@{}}\toprule
Open dependency & Discriminating evidence & Decision changed & Claim still not proved\\\midrule
Physical manipulation & Direct field/light/noise calibration and sham equivalence. & Whether a cue-conflict test is interpretable. & Animal uses the cue.\\
Cue dependence & Predicted contrast plus restoration and complete individual distribution. & Whether to pursue the cue under this context. & A map, route benefit, or receptor.\\
Mechanism bridge & Target-specific perturbation, rescue, sensory-function and health controls. & Whether the candidate pathway remains plausible. & In vivo complete receptor mechanism.\\
Ecological consequence & Track provenance, weather/current adjustment, missingness sensitivity, route prediction. & Whether field deployment is warranted. & Species-wide migration explanation.\\
Cross-context extension & Independent site, season, or taxon using locked protocol. & Whether the boundary expands or contracts. & A cue common to all migratory animals.\\\bottomrule
\end{tabular}
\end{table*}

The reproducibility package should contain protocol version, species and route boundary, cue calibration records, sensor and device metadata, animal welfare events, release and track provenance, all inclusion and exclusion rules, raw and cleaned bearing data, analysis code, null tests, negative results, and a statement of the strongest conclusion permitted. This paper contains no empirical package because it reports no study outcome.

\section{Cross-Taxon Predictions and Comparative Design}\label{sec:comparative}
The framework must resist a second kind of overgeneralization: treating a useful bird experiment as a generic animal-navigation experiment. A songbird's light-dependent retinal candidate, an insect's time-compensated sun compass, a sea turtle's open-ocean orientation, and an electroreceptive fish's induction-sensitive environment place different constraints on the intervention. Cross-taxon work is strongest when it predicts a contrast from anatomy, habitat, optical environment, locomotion, or known sensory repertoire before observing the outcome.

For example, a radical-pair candidate predicts sensitivity to the joint light regime and field geometry, whereas an induction account predicts dependence on movement, conductivity, and electrosensory state. A landmark or olfactory hypothesis predicts stronger effects near familiar regions or on routes where the corresponding information is present. These predictions can all be wrong, but they make an experiment discriminating. A cue is not shown to be unimportant merely because its prediction fails in a condition in which it was never physically available or ecologically reliable.

Comparative tests should first preserve a common inferential core: physical cue calibration, sham and restoration conditions, individual-level responses, welfare monitoring, and transparent attrition. The species-specific layer then changes the intervention. A magnetic coil arrangement appropriate to a songbird may be irrelevant or unsafe for a marine animal; a virtual magnetic displacement may be more interpretable than physical relocation for a turtle; and a clock-shift design can probe sun-compass computation in insects without presupposing an avian receptor.

The appropriate result of comparative work is a \emph{mechanism map}, not a species ranking. The map identifies which cues have evidence at what level and under which context. It allows common functions to emerge without demanding identical receptors. It also makes genuine divergence informative: two animals may both follow a migratory bearing while deriving it from different sensory pathways or using the same cue for different computations.

\begin{table*}[!t]
\caption{Comparative predictions that keep function separate from receptor identity.}\label{tab:comparative}\centering\footnotesize
\begin{tabular}{@{}L{.18\textwidth}L{.22\textwidth}L{.28\textwidth}L{.24\textwidth}@{}}\toprule
System boundary & Leading discriminating question & Context-appropriate manipulation and control & Result that would narrow, not universalize, the claim\\\midrule
Nocturnal migratory songbird & Does a light-dependent magnetic cue contribute to a compass response? & Field rotation/displacement across controlled spectra, sham coils, restoration, visual-function controls. & A spectral and field-dependent orientation effect for the stated population and season.\\
Monarch or other diurnal insect & Does circadian state calibrate solar or polarized-light direction? & Clock shift, sky cue rotation, antenna or sensory intervention where ethically justified, matched handling. & A time-dependent celestial-compass contribution under the tested sky condition.\\
Sea turtle or other marine migrant & Does field information contribute to open-sea orientation or displacement correction? & Virtual field displacement with oceanographic, release, and tag-quality provenance. & A field-context route correction hypothesis for a defined life stage and release region.\\
Electroreceptive aquatic animal & Is an induction-like signal available during movement through conductive water? & Joint flow, conductivity, field, and electroreception controls. & Conditional dependence on the predicted physical interaction.\\
Familiar-route terrestrial migrant & Do landscape or odour cues calibrate a regional map? & Familiar/unfamiliar release and cue masking with welfare and sensory controls. & A local or regional cue contribution, not a global navigation theory.\\\bottomrule
\end{tabular}
\end{table*}

\section{Data Architecture, Welfare, and Stopping Rules}\label{sec:governance}
Navigation studies often combine rich biological signals with fragile observations. A high-resolution track can be missing exactly when weather is poor, a magnetic manipulation can be applied unevenly across a release block, and a behavioural score can be unknowingly influenced by an observer who knows the arm. MAGNAV treats these as scientific risks rather than later data-cleaning inconveniences. Every stage has a physical-provenance record, an animal-provenance record, and an analytical-provenance record.

Physical provenance includes the three-dimensional magnetic vector, local anomaly and time variation, coil calibration, RF noise, illumination spectrum and irradiance, sky conditions, weather, wind or current, and release geometry. Animal provenance includes capture or observation time, condition assessment, handling duration, device characteristics, release time, social context, and ethical withdrawal. Analytical provenance includes the preregistered estimator, code version, all raw and derived files, track-quality rule, blind status, and the reason for each exclusion. The record is useful even when the experiment is negative, because it distinguishes an absent behavioural response from an unverified cue manipulation.

Welfare is noncompensatory. A study cannot justify a mechanistic claim by producing a large orientation contrast if device burden, stress, or abnormal sensory function make the contrast uninterpretable. The protocol therefore defines stop rules for individual distress, tag or device malfunction, unexpected mortality, adverse weather, unacceptable release conditions, and loss of calibrated manipulation. These rules must be applied independently of whether preliminary responses resemble the expected direction.

Stopping rules also protect inferential integrity. If the sham field drifts, the light spectrum changes outside tolerance, a route prediction was not fixed, or tracking loss crosses a predeclared threshold, the study pauses or downgrades its claim rather than continuing to accumulate an opaque average. The next action can be recalibration, a narrower sensory claim, or a redesigned replication. It cannot be retrospective relabelling of the problem as biological variability.

\begin{table*}[!t]
\caption{Operational records and stopping rules for a MAGNAV run.}\label{tab:governance}\centering\footnotesize
\begin{tabular}{@{}L{.20\textwidth}L{.27\textwidth}L{.25\textwidth}L{.20\textwidth}@{}}\toprule
Domain & Required evidence ledger & Pause or stop trigger & Claim consequence\\\midrule
Physical cue delivery & Vector, spectrum, noise, temperature, hardware, timing, and sham equivalence. & Calibration outside tolerance or unresolvable arm difference. & No cue-specific analysis for affected block.\\
Animal welfare and sensory function & Condition checks, handling, device burden, recovery, adverse events, sensory control. & Distress, abnormal impairment, unexpected mortality, or protocol breach. & Withdraw animal; do not code as navigational failure.\\
Tracking and data availability & Positional uncertainty, battery/duty cycle, contact loss, environmental coverage, time at risk. & Selective loss beyond declared sensitivity bound. & Restrict ecological estimand or suspend inference.\\
Analysis integrity & Frozen hypotheses, randomization, blinding, code, exclusions, and null tests. & Broken blind, unlogged protocol deviation, or post hoc primary outcome change. & Label exploratory or repeat with locked protocol.\\
Communication & Claim level, negative results, counterexamples, context boundary, mechanism status. & Language exceeds the strongest passed gate. & Revise statement before dissemination.\\\bottomrule
\end{tabular}
\end{table*}

\section{Interface Failure Modes and Threats to Validity}\label{sec:threats}
The most common error occurs at the interface between measurement and interpretation. A changed bearing can be an orientation response without being a navigational map. A response to a virtual magnetic displacement may show access to field information without proving that an animal computes a geographic coordinate. A molecular magnetic effect can constrain mechanistic theories without identifying a sensory receptor. A route shift can be ecologically important without being caused by the manipulated cue.

\begin{table*}[!t]
\caption{Interface failure matrix for migration-navigation inference.}\label{tab:interfaces}\centering\footnotesize
\begin{tabular}{@{}L{.18\textwidth}L{.25\textwidth}L{.29\textwidth}L{.20\textwidth}@{}}\toprule
Interface & Example failure & Evidence needed before a stronger claim & Consequence for interpretation\\\midrule
Cue to bearing & Apparatus, handling, or altered illumination changes preferred direction. & Sham, physical calibration, blinded scoring, restoration, and retained-cue ledger. & Observation is not cue-specific.\\
Bearing to map & Animal follows a compass but cannot correct a displacement. & Virtual-displacement or release design with position-dependent predictions. & Compass result is not a map result.\\
Molecule to behaviour & Candidate protein responds in vitro but is not used in sensory tissue during migration. & Localization, state-relevant expression, perturbation/rescue, sensory-function controls. & Molecular result is not a receptor claim.\\
Lab to field & Cue use in an arena disappears under wind, landscape, or social context. & Track-quality, environmental provenance, route prediction, and ecological replication. & Arena effect has limited transferability.\\
One taxon to many & A mechanism works in a songbird but is generalized to insects or marine animals. & Independent lineage-specific sensory and ecological tests. & Claim remains within the tested boundary.\\\bottomrule
\end{tabular}
\end{table*}

Selection is another threat. A target cue can appear decisive if animals are chosen after a promising orientation assay, if a field season has unusually clear skies, or if track loss selectively removes animals in difficult conditions. MAGNAV requires the denominator of attempted measurements, device failures, welfare withdrawals, and environmental exclusions. It also requires the mapping between a claimed cue and the measured physical signal. Without that mapping, a label such as ``magnetic condition'' can conceal a mixture of field, light, hardware, and handling differences.

Transportability needs the same discipline. An effect seen in one latitude, season, or stage can be genuine and still not generalize. The scientifically useful outcome is often a narrower boundary: for example, a cue matters for inexperienced juveniles under overcast conditions but not for adults on familiar routes. Restriction is a discovery about the computation; it should not be treated as an inconvenient deviation from a universal narrative.

\section{Discussion and Conclusion}\label{sec:discussion}
Migratory animals do not need one magical sense to navigate. The strongest current account is that animals integrate cues across spatial and temporal scales, switching or weighting them as their availability and reliability change. Celestial patterns, visual landscapes, odours, hydrodynamics, biological clocks, and magnetic information can each be valuable, but their roles depend on the organism and task. The 2018 long-distance-navigation review makes this integration explicit \cite{mouritsen2018}; the 2021 CRY4 work adds important molecular evidence for one radical-pair candidate in one migratory bird \cite{xu2021}.

MAGNAV makes this a constructive research programme. It asks a candidate explanation to pass the tests required by its breadth: physical manipulation integrity; conditional cue-dependent orientation; mechanism-to-behaviour connection; ecological route contribution; and independent replication. It treats a negative or context-specific result as information that can refine the cue-dependence ledger. It does not average away species differences or let a positive molecular assay stand in for a free-ranging migratory calculation.

This four-stage design-only workflow concludes that the question ``How do migratory animals know where they are going?'' should be answered through bounded, multiscale evidence rather than a single universal mechanism. Future studies can demonstrate a directional cue, then a context-dependent weight, then a receptor or circuit contribution, then a field consequence for a stated population. Each step can be important. No finite set of early results justifies the claim that all migratory animals use the same compass, map, receptor, or navigation algorithm.

\section{Research Priorities}\label{sec:priorities}
The highest-value next measurements are those that discriminate between competing cue-integration accounts, not merely those that record another bearing. First, field and laboratory studies should report the physical cue environment with enough resolution to establish what the animal could have sensed: three-dimensional magnetic vector and variability, light spectrum and irradiance, RF noise, sky condition, wind or current, and landscape or odour availability. This converts a nominal ``magnetic'' or ``celestial'' manipulation into an auditable exposure definition.

Second, mechanism studies should be linked to the behavioural context for which they are proposed. Molecular work on cryptochromes, neural recordings, visual or trigeminal interventions, and magnetite investigations are most informative when their expression, localization, timing, and control phenotypes are measured in the same state in which the relevant orientation behaviour is tested. A mechanistic experiment that excludes competing sensory impairment can substantially narrow theory even before a free-ranging route test is feasible.

Third, comparative programmes should coordinate standards rather than force a common receptor narrative. Shared metadata, calibrated manipulations, complete negative results, and open cue-dependence ledgers allow genuine commonalities and lineage-specific differences to emerge. In this setting, an unsuccessful prediction can be valuable: it may identify the ecological condition, developmental phase, or sensory combination in which a cue is absent, downweighted, or replaced. MAGNAV treats that restriction as progress toward a mechanism map.

\appendices
\section{Minimal MAGNAV Run Manifest}
Each future run records species and population; migratory phase and route segment; individual age/experience and social context; cue definitions and calibrated physical values; retained sensory information; manipulation, sham, and restoration specifications; welfare and stop rules; sampling denominator; field, light, sky, weather, wind, water-current, and RF records; device characteristics; bearing and tracking endpoints; missingness, exclusion, and censoring ledger; predeclared analysis; null and negative controls; and the maximum permitted claim.

\section{Protocol Checklist}
Before data collection, freeze the population boundary, cue conflict, expected response, field and light calibration, sham, restoration, outcome definitions, quality thresholds, route prediction, and ethical stop rules. During collection, preserve physical measurements, handling and device records, failed tracks, sensory-function checks, individual-level data, and all negative results. Before communication, label evidence as cue availability, conditional orientation, mechanism bridge, ecological route contribution, or replication; do not convert a finite context-specific result into a universal explanation of migration.

\newpage
\begin{thebibliography}{10}
\bibitem{mouritsen2018}H. Mouritsen, ``Long-distance navigation and magnetoreception in migratory animals,'' \emph{Nature}, vol. 558, pp. 50--59, 2018, doi: 10.1038/s41586-018-0176-1.
\bibitem{xu2021}J. Xu \emph{et al.}, ``Magnetic sensitivity of cryptochrome 4 from a migratory songbird,'' \emph{Nature}, vol. 594, pp. 535--540, 2021, doi: 10.1038/s41586-021-03618-9.
\bibitem{ritz2000}T. Ritz, S. Adem, and K. Schulten, ``A model for photoreceptor-based magnetoreception in birds,'' \emph{Biophysical Journal}, vol. 78, no. 2, pp. 707--718, 2000, doi: 10.1016/S0006-3495(00)76629-X.
\bibitem{hore2016}P. J. Hore and H. Mouritsen, ``The radical-pair mechanism of magnetoreception,'' \emph{Annual Review of Biophysics}, vol. 45, pp. 299--344, 2016, doi: 10.1146/annurev-biophys-032116-094545.
\bibitem{wiltschko2005}R. Wiltschko and W. Wiltschko, ``Magnetic orientation and magnetoreception in birds and other animals,'' \emph{Journal of Comparative Physiology A}, vol. 191, pp. 675--693, 2005, doi: 10.1007/s00359-005-0627-7.
\bibitem{zapka2009}M. Zapka \emph{et al.}, ``Visual but not trigeminal mediation of magnetic compass information in a migratory bird,'' \emph{Nature}, vol. 461, pp. 1274--1277, 2009, doi: 10.1038/nature08528.
\bibitem{merlin2009}C. Merlin, R. J. Gegear, and S. M. Reppert, ``Antennae are circadian clocks that orchestrate sun compass navigation in monarch butterflies,'' \emph{Science}, vol. 325, no. 5948, pp. 1700--1704, 2009, doi: 10.1126/science.1176228.
\bibitem{muheim2006}R. Muheim, J. B. Phillips, and S. Akesson, ``Polarized light cues under overcast skies: the role of avian skylight navigation,'' \emph{Journal of Experimental Biology}, vol. 209, pp. 243--259, 2006, doi: 10.1242/jeb.02018.
\bibitem{lohmann1996}K. J. Lohmann and C. M. F. Lohmann, ``Orientation and open-sea navigation in sea turtles,'' \emph{Journal of Experimental Biology}, vol. 199, pp. 73--81, 1996.
\end{thebibliography}
\end{document}
